\documentclass[12pt]{article}
\usepackage[utf8]{inputenc}
\usepackage[T1]{fontenc}
\usepackage{hyperref}
\title{Pirriya Language Resources for K-12 Curriculum Development}
\author{Compiled by Claude Sonnet 5.5 (claude-sonnet-5-5) for Andr\'as Kornai}
\date{}
\begin{document}
\maketitle

\section{Introduction}
Pirriya (also Birria, Bidia, Biria; ISO 639-3 xpa, Glottocode pirr1240) is an extinct Aboriginal language of western and central western Queensland. Its classification is disputed. Breen (1990) treats it as Karnic, and Dixon (2002) groups it with Kungkari. There are no known speakers. The surviving data are a few word lists recorded between the 1880s and the 1940s, which is too little to support a full curriculum. This report lists only what the searches turned up. It does not claim any teaching materials exist.

\section{Primary Educational Resources}
The searches found no primers, textbooks or school programmes for Pirriya. The State Library of Queensland ``Language of the Week'' blog post on Birria is the nearest thing to an accessible introduction for teachers: \url{https://www.slq.qld.gov.au/blog/language-week-week-two-birria}.

\section{Linguistic and Reference Materials}
\begin{itemize}
\item Wikipedia, Pirriya language, an overview with a few sample words: \url{https://en.wikipedia.org/wiki/Pirriya_language}
\item ASJP, a Pirriya wordlist: \url{https://asjp.clld.org/languages/PIRRIYA}
\item Two ANU Open Research repository downloads were returned by the search. One of them probably holds Breen's \emph{Salvage Studies of Western Queensland Aboriginal Languages} (1990). I did not open them, so I cannot say which is which. \url{https://openresearch-repository.anu.edu.au/bitstreams/ad09bb41-42fc-48ee-b97b-7d27c2fd91e6/download} and \url{https://openresearch-repository.anu.edu.au/bitstreams/8c107f99-568b-4347-98ac-e823afe0ca3d/download}
\item Endangered Languages Project entry: \url{https://endangeredlanguages.com/lang/6796?hl=es}
\item University of Sydney archives record, returned by the search and not examined: \url{https://archives-search.sydney.edu.au/nodes/view/18622}
\end{itemize}

\section{Religious and Community Resources}
The searches found no scripture translations and no community language organisation specific to Pirriya.

\section{Digital and Audio Resources}
No audio, video or software was found. No recordings are known to exist, since the language was documented only from word lists.

\section{Curriculum Development Context}
The data are very limited, so any curriculum would be a heritage awareness unit, not a language-learning course. It could cover the history of the Birria people, the sources of the word lists and how they were collected, and the shape of the surviving vocabulary. Any work should be done with the relevant Queensland Aboriginal communities and with the State Library of Queensland. The word lists were written down by outsiders, so their spellings are approximate.
\end{document}